\section{Evaluation: Five Failure Scenarios}
\label{sec:evaluation}

We evaluate $\eps$ across five scenarios spanning four distinct failure
classes. All five scenarios use GPT-4o-2024-08-06 via the OpenAI API with
\texttt{logprobs=True, top\_logprobs=5}. The system prompt specifies
naming conventions (concise conventional parameter names, direct verb-noun
function names) to suppress Type~1 cosmetic uncertainty before aggregation.
Each scenario is implemented as a standalone script and is included in the
repository. Multi-model validation extending these results to two additional
OpenAI models appears in Section~\ref{sec:eval-multimodel}.

\subsection{Scenario A: Deprecated API}
\label{sec:eval-a}

\paragraph{Setup.}
The Stripe payments API underwent a major deprecation between v6 and v7:
\texttt{stripe.Charge.create()} was removed in favour of
\texttt{stripe.PaymentIntent.create()}, which requires a two-step
intent-then-confirm flow. Both APIs are heavily represented in the model's
training data, as the deprecation occurred during a period of active
online documentation.

\paragraph{Prompt.}
\textit{``Write a Python function to charge a customer's card for \$50 using
the Stripe Python SDK in a FastAPI application. The customer\_id is passed
in as a parameter.''} Context: Python 3.11, Stripe SDK installed.

\paragraph{Result: \textsc{paused}, $\eps = 0.914$.}
The highest-entropy token in the generated function is the method selector
immediately following the \texttt{stripe} module reference: the token
\texttt{.Ch} (beginning of \texttt{Charge}) appears with probability 0.52,
while \texttt{.Payment} (beginning of \texttt{PaymentIntent}) appears with
probability 0.47 --- a near-50/50 split on a production-critical decision.
Secondary flags include route URL naming (\texttt{/charge\_customer} vs.\
\texttt{/charge-customer}, $\eps = 0.75$) and return value design
(\texttt{payment\_intent.id} vs.\ the full object, $\eps = 0.48$).

\paragraph{Finding.}
The model's uncertainty faithfully reflects the state of its training data:
both API paths are real, both are documented, and neither is definitively
preferred by the available evidence. A developer reviewing this \textsc{paused}
output knows to verify which API version their installed Stripe SDK supports
before committing.

\subsection{Scenario B: Confident Wrongness}
\label{sec:eval-b}

\paragraph{Setup.}
The OpenAI Python SDK introduced a breaking API change between v0 and v1.
The v0 interface (\texttt{openai.ChatCompletion.create(...)}) raises an
\texttt{AttributeError} on any SDK version above 1.0. GPT-4o-2024-08-06
was trained predominantly on v0 documentation and generates v0 syntax
with high confidence on a neutral prompt.

\paragraph{Prompt.}
\textit{``Write a Python function that calls GPT-4 with a user message and
returns the response text.''} Context: Python 3.11, openai package installed.

\paragraph{Result: \textsc{paused}, $\eps = 0.551$.}
The model generated v0 syntax --- \texttt{openai\allowbreak.ChatCompletion\allowbreak.create(\ldots)}
with \texttt{response\allowbreak[`choices']\allowbreak[0]\allowbreak[`message']\allowbreak[`content']}
access --- with high confidence. The \textsc{paused} status was triggered not by the API
choice itself but by uncertainty in the function's parameter design: the
model was uncertain whether to accept a raw string
(\texttt{user\_message}), a pre-formatted dict (\texttt{message}), or a
full messages list (\texttt{messages}), and whether to expose the model
name as a parameter.

\paragraph{Finding.}
This scenario demonstrates $\eps$'s honest boundary. The system cannot
detect a decision the model makes confidently and incorrectly: it reports
``the model was certain about \texttt{openai.ChatCompletion.create}'' ---
which is true. What $\eps$ provides is \emph{review pressure}: the
\textsc{paused} status prompts the developer to examine the generated
function, during which they encounter the deprecated v0 call. The indirection
is not a failure; it is the mechanism. A tool that drove review of every
uncertain decision would surface this bug whenever the model was uncertain
about anything in the same generation.

\subsection{Scenario C: Syntax Split}
\label{sec:eval-c}

\paragraph{Setup.}
SQLAlchemy 1.4 introduced a compatibility layer and SQLAlchemy 2.0
completed a full ORM query API migration. The 1.x session query interface
(\texttt{db.query(Model).filter(...)}) and the 2.0 execute interface
(\texttt{db.execute(select(Model)).scalars().all()}) coexist in the model's
training data due to the long overlap period.

\paragraph{Prompt.}
\textit{``Write a Python function that queries all active users from a
SQLAlchemy database session and returns them as a list. The codebase has
been around for a few years.''} Context: Python 3.11, SQLAlchemy installed.

\paragraph{Result: \textsc{flagged}, $\eps = 0.435$.}
The model chose 1.x syntax (\texttt{db.query(User).filter(User.active ==
True).all()}). The primary flag was on the import path: the model was
uncertain whether to import \texttt{User} from a local \texttt{models}
module or construct the select via \texttt{sqlalchemy.orm}. The query
method decision itself fell below the flag threshold, consistent with the
model having a moderate preference for 1.x.

\paragraph{Finding.}
\textsc{flagged} rather than \textsc{paused} correctly reflects the
situation: the model has a preference, but it is not certain. A developer
reviewing this output knows to verify which SQLAlchemy version is pinned
in \texttt{requirements.txt}. If the project targets SQLAlchemy 2.0,
the generated query will emit a deprecation warning and will eventually
break.

\subsection{Scenario D: Compatibility Mismatch}
\label{sec:eval-d}

\paragraph{Setup.}
This scenario introduces a qualitatively different failure class: not a
single uncertain decision but two \emph{locally reasonable} decisions that
combine into a silent bug. In a FastAPI application, using
\texttt{async def} for an endpoint handler and \texttt{requests.get()} for
the HTTP call is individually defensible --- but the combination blocks
the event loop under concurrent load, degrading the entire application to
synchronous throughput.

\paragraph{Prompt.}
\textit{``Write a FastAPI endpoint that accepts a city name as a query
parameter, fetches the current weather for that city from an external
weather API, and returns the result as JSON.''} Context: Python 3.11,
FastAPI installed, requests library used throughout the existing codebase.

\paragraph{Result: \textsc{paused}, $\eps = 0.812$.}
The model generated \texttt{async def get\_weather(...)} with
\texttt{requests.get(url)} inside the body --- the mismatch. Notably,
neither decision individually produced a high-$\eps$ token: the
\texttt{async} keyword token showed $\eps = 0.26$ (the model had a mild
preference for async given the FastAPI context), and the \texttt{requests}
import showed $\eps = 0.31$ (the model correctly reflected that requests
was the established library in the described codebase). The elevated
file-level $\eps$ came from uncertainty in the URL construction and API key
handling, not from the compatibility decision itself. A mismatch-detection
pass --- comparing the presence of \texttt{async def} against blocking
HTTP client calls in the function body --- surfaced the bug explicitly in
the output.

\paragraph{Finding.}
This is the hardest class of code generation failure to detect: two
individually confident decisions that are jointly incompatible.
$\eps$ does not directly flag the compatibility issue --- the model was
not uncertain about either component. What $\eps$ provides is the
\textsc{paused} status that prompts review, and the mismatch-detection
post-processor that makes the compatibility violation explicit when it
exists. The failure class argues for extending $\eps$ with cross-decision
consistency checks as a future contribution.

\subsection{Scenario E: Multi-Function Module}
\label{sec:eval-e}

\paragraph{Setup.}
Scenarios A--D each isolate a single consequential decision. Scenario~E
presents four simultaneous decisions that are also interdependent: password
hashing library (bcrypt / argon2 / passlib / hashlib), JWT library (PyJWT
/ python-jose / authlib), SQLAlchemy style (1.x / 2.0), and datetime
handling (\texttt{datetime.utcnow()} deprecated vs.\ timezone-aware
\texttt{datetime.now(timezone.utc)}). These decisions must be made
consistently across six functions: \textsc{register\_user}, \textsc{login},
\textsc{verify\_token}, \textsc{refresh\_token},
\textsc{request\_password\_reset}, and \textsc{reset\_password}.

\paragraph{Prompt.}
A complete Python authentication module with the above six functions,
specifying their contracts but not their implementation. Context: Python
3.11, FastAPI backend, PostgreSQL via SQLAlchemy, no existing authentication
library committed.

\paragraph{Result: \textsc{paused} on all 16 runs, $\eps \in [0.69, 0.92]$.}
Across 16 independent runs, the model chose different libraries: passlib
(2 runs), bcrypt (1 run), PyJWT (2 runs), python-jose (1 run), and mixed
combinations in the remainder. The model's cross-run inconsistency reflects
genuine uncertainty about the right choice --- but any single run \emph{looks}
confident, because the model commits to one library in the first function
and uses it consistently within that generation.

A consistent intra-run pattern emerged: functions generated early in the
context (\textsc{register\_user}, \textsc{login}) showed $\eps = 0.40$--$0.58$
at the library import decision; functions generated later
(\textsc{verify\_token}, \textsc{reset\_password}) frequently showed
$\eps = 0.000$ at those same positions, because the model had already
committed to a library and was executing within its established convention.
This \emph{uncertainty collapse} is itself a meaningful signal: the
high-risk decision window is the first function of each library domain.

Per-function $\eps$ breakdown and cross-function consistency analysis ---
flagging cases where the model used different JWT libraries in
\textsc{login} and \textsc{verify\_token} --- are provided by the demo
output.

\paragraph{Finding.}
The multi-function setting reveals both the power and the limit of $\eps$.
Within a single generation, it correctly identifies the first-function
decision point as the highest-risk location. Across generations, the
variation in library choice confirms that the model's uncertainty is real.

\subsection{Summary}

Table~\ref{tab:scenarios} and Figure~\ref{fig:scenarios} summarise the
five scenarios. The four failure classes --- deprecated API, confident wrongness, syntax
split, and compatibility mismatch --- are qualitatively distinct.
The first and third are detectable directly by $\eps$. The second is
detectable indirectly, through the review pressure $\eps$ creates. The
fourth requires augmenting $\eps$ with cross-decision consistency checks.

\begin{table}[h]
\centering
\caption{Scenario summary. All runs on GPT-4o-2024-08-06.
         ``Indirect'' detection (B) means $\eps$ drives review that surfaces
         the bug, not that $\eps$ directly flags the erroneous token.}
\label{tab:scenarios}
\begin{tabular}{llccc}
\toprule
\textbf{Scenario} & \textbf{Failure class} & \textbf{$\eps$}
                  & \textbf{Status} & \textbf{Detection} \\
\midrule
A --- Stripe deprecation    & Deprecated API      & 0.914      & \textsc{paused}  & Direct \\
B --- OpenAI SDK v0         & Confident wrongness & 0.551      & \textsc{paused}  & Indirect \\
C --- SQLAlchemy style      & Syntax split        & 0.435      & \textsc{flagged} & Direct \\
D --- FastAPI async/sync    & Compatibility       & 0.812      & \textsc{paused}  & Augmented \\
E --- Auth module (16 runs) & Multi-decision      & 0.69--0.92 & \textsc{paused}  & Direct (all runs) \\
\bottomrule
\end{tabular}
\end{table}

\begin{figure}[H]
\centering
\includegraphics[width=0.82\textwidth]{fig3_scenarios.png}
\caption{Five evaluation scenarios with $\varepsilon$ score and failure class.
         Orange dashed line: soft threshold (0.30, \textsc{flagged}).
         Red dashed line: hard threshold (0.65, \textsc{paused}).
         Scenario~E bar shows the 16-run mean; the annotation gives the full
         range $[0.69, 0.92]$.}
\label{fig:scenarios}
\end{figure}

\subsection{Multi-Model Validation}
\label{sec:eval-multimodel}

To assess whether the $\eps$ signal is specific to GPT-4o or holds across
models, we re-ran the 30-prompt benchmark on two additional OpenAI models:
\texttt{gpt-4o-mini-2024-07-18} and \texttt{gpt-4-turbo-2024-04-09}.
All settings were identical to the five-scenario runs (same prompts,
same system prompt, \texttt{logprobs=True, top\_logprobs=5}).

\begin{table}[h]
\centering
\caption{Cross-model detection results on the 30-prompt benchmark
         (20 HIGH prompts, 10 LOW prompts). Detection rate is on HIGH prompts;
         FP rate is on LOW prompts.}
\label{tab:multimodel}
\begin{tabular}{lccc}
\toprule
\textbf{Model} & \textbf{HIGH detection} & \textbf{LOW FP rate} & \textbf{Note} \\
\midrule
GPT-4o-2024-08-06       & \textbf{85\%} (17/20) & 50\% (5/10) & Best calibrated \\
GPT-4o-mini-2024-07-18  & 70\% (14/20)          & 60\% (6/10) & Lower det.\ + higher FP \\
GPT-4-turbo-2024-04-09  & 70\% (14/20)          & 60\% (6/10) & Same rate; different missed set \\
\bottomrule
\end{tabular}
\end{table}

\begin{figure}[H]
\centering
\includegraphics[width=0.82\textwidth]{fig4_multimodel.png}
\caption{$\varepsilon$ detection rate (HIGH prompts) and false positive rate
         (LOW prompts) across three OpenAI models on the 30-prompt benchmark.
         Detection correlates with model calibration quality (GPT-4o highest
         at 85\%; two additional models at 70\%).}
\label{fig:multimodel}
\end{figure}

The $\eps$ signal holds across all three models. GPT-4o achieves the highest
detection rate, which we attribute to better logprob calibration: when
GPT-4o is genuinely uncertain about an API version decision, that uncertainty
is more accurately reflected in the token probability distribution than in
smaller or older models. This pattern is consistent with the hypothesis that
logprob calibration quality affects $\eps$ signal strength --- GPT-4o's
better-calibrated token distributions may reflect genuine uncertainty more
accurately than older or smaller models. Disentangling calibration from
capability, training-data recency, and sampling variance would require
broader model coverage across vendors; we leave this as future work.

\paragraph{Consistent misses across all models.}
Three prompts produced $\eps_\text{file} = 0.000$ across all models and all
runs: current UTC datetime retrieval, HTTP GET to JSON parsing, and (for
GPT-4-turbo) PostgreSQL query construction. These are domains where every
model has converged on a strong training-data prior: the API choice is
effectively resolved. This is consistent with the $\eps$ story, not a
counter-example to it --- these domains have resolved their training-data
split, so there is genuinely no uncertainty to detect. We term these
\emph{confidently-resolved domains}; their limitations are discussed in
Section~\ref{sec:discussion}.
